% !PRINTED 389
\weberpage{389}
%<<CONT
\gesperrt{exclude numbers which are not relatively prime to $S$ in the numerator or in the denominator.} (In the case of the orders, the prime factors of the conductor $Q$ must be contained in $S$.)

We now form from $M$ a divisor $A$, into which we take all numbers $a$ that satisfy a congruence

\weq{2}{N(\omega) \equiv a \pmod{m}}

that is, for which there exists a number $\omega$ in $O$ whose norm is congruent to $a$.

1. \gesperrt{These numbers are called norm residues for the modulus $m$.}

From the definition of the norm residues it follows that the norm residues form a group with respect to multiplication and division.

It follows further that \gesperrt{every square and every quadratic residue (modulo $m$) is at the same time a norm residue, and hence that every norm non-residue is at the same time a quadratic non-residue.} From the group nature of $A$ it follows:

\gesperrt{The product of a norm residue and a norm residue is a norm residue.}

\gesperrt{The product of a norm residue and a norm non-residue is a norm non-residue.}

The group of norm residues always has a finite index $(Z, A)$, for it unites within itself whole \gesperrt{classes of numbers} with respect to the modulus $m$, and the number of these is finite.

\S\,107. \gesperrt{Composition of the norm residue groups.}

We wish to consider two moduli $m_1, m_2$ that are relatively prime to one another, and we denote the groups of norm residues of $m_1$ and $m_2$ by $A_1$ and $A_2$. If we set for the moment

\[ (Z, A_1) = \mu \]

and decompose

\weq{1}{Z = a_1 A_1 + a_2 A_1 + \cdots + a_\mu A_1,}

then we may replace the $a_1, a_2, \ldots, a_\mu$ by arbitrary numbers congruent to them with respect to the modulus $m_1$. Since $m_1$ and $m_2$ are assumed to be relatively prime, rational $x_i$ can be determined, whatever the rational numbers $c_i$ may be, from the congruences

\[ a_i + x_i m_1 \equiv c_i \pmod{m_2} \]

and one can therefore choose the $a_i$ so that they are contained in $A_2$. One can therefore represent every number $z$ in $Z$ as a product
%CONT>>
